
\subsubsection*{5.2.5}
\phantomsection
\label{sga3:VIIB:5.2.5}

It remains to consider the case in which \(B\) has infinite
height and the projection \(I\to\omega\) admits a section
\(\sigma\) compatible with the filtrations of \(\omega\) and
\(I\). Consider the morphism induced by \(\sigma\):
\[
  \phi:
  B'
  =
  k[[\omega]]
  \big/
  \bigl((x^{p^n})\bigr)_{x\in\omega_n}
  \longrightarrow B.
\]
It is enough to prove that, for every \(r\in\overline{\mathbb N}\),
the map
\[
  \phi_r:B'_r\longrightarrow B_r
\]
induced by \(\phi\) is an isomorphism.\footnote{\textit{Editors' note.} %
\label{sga3:VIIB:nde:162}Indeed, \(B'\), respectively \(B\), is
the inverse limit of the \(B'_r\), respectively the \(B_r\).
Moreover, the editors modified the original in what follows in
order to take account of the addition made in \SGARefLink{sga3:VIIB:5.2.3}{5.2.3}.}

For every \(r\in\mathbb N^*\), let \(G_r\) be the kernel of the
iterated Frobenius morphism
\[
  G\longrightarrow G^{(p^r)},
\]
and let \(HG_r\) be the formal subgroup of \(G\) that is the
inverse image of \(H^{(p^r)}\subset G^{(p^r)}\). Thus \(HG_r\)
is defined by the closed ideal generated by the \(p^r\)-th
powers of the elements of \(AI\), namely
\[
  A\check J_r,
  \qquad
  \check J_r=\{\,x^{p^r}\mid x\in I\,\}.
\]
Exactly as in \SGARefLink{sga3:VIIB:5.2.3}{5.2.3}, it follows that
\[
  B_r=B/B\check J_r
\]
is the affine algebra of \(HG_r/H\); in particular, it has
height at most \(r\).

Let
\[
  \mathfrak m(r)=I/B\check J_r
\]
be the augmentation ideal of \(B_r\). The canonical projection
\(B\to B_r\) induces an isomorphism
\[
  \omega=I/I^2
  \xrightarrow{\ \sim\ }
  \omega(r)
  =
  \mathfrak m(r)/\mathfrak m(r)^2,
\]
which we use to identify the two spaces. Let \(\omega(r)_n\) be
the image in \(\omega(r)\) of the closed ideal
\[
  \mathfrak m(r)_n
  =
  \{\,y\in\mathfrak m(r)\mid y^{p^n}=0\,\}.
\]
It is also the image in \(\omega\) of the closed ideal
\[
  I(r)_n
  =
  \{\,x\in I\mid x^{p^n}\in B\check J_r\,\}.
\]
\footnote{\label{sga3:VIIB:nde:163}The editors expanded the
original in what follows.}

Clearly,
\[
  \omega(r)_n=\omega
  \qquad\text{if }n\geq r.
\]
We shall prove that
\[
  \omega(r)_n=\omega_n
  \qquad\text{if }n<r.
\]
For every \(r,n\), the sequence
\[
  0
  \longrightarrow
  I(r)_n\cap I^2
  \longrightarrow
  I(r)_n
  \longrightarrow
  \omega(r)_n
  \longrightarrow0
\]
is exact. For fixed \(n\),
\[
  \bigcap_r I(r)_n=I_n,
\]
because \(\bigcap_rB\check J_r=0\). Filtered inverse limits are
exact in \(\operatorname{PC}(k)\) \textup{(cf.\ \SGARefLink{sga3:VIIB:0.2}{0.2})}, so for
every \(n\),
\[
\tag{*}\label{sga3:vii:remark:5-2-4-a:equation:star}
  \omega_n=\bigcap_r\omega(r)_n.
\]
On the other hand, \SGARefLink{sga3:VIIB:5.2.4.A}{Remark~5.2.4.A} gives
\[
  \omega(r)_n=\omega(r+1)_n
  \qquad\text{if }n<r.
\]
Together with \textup{(*)}, this implies
\[
  \omega(r)_n=\omega_n
  \qquad\text{if }n<r.
\]

Consequently, the vector space \(\omega\), filtered by the
subspaces \(\bigl(\omega(r)_n\bigr)_{n\geq0}\), is precisely
\({}^{r}\omega\) in the notation of \SGARefLink{sga3:VIIB:5.2:theorem}{5.2}. The composite
\[
  \sigma(r):
  \omega\xrightarrow{\ \sigma\ }I
  \longrightarrow\mathfrak m(r)
\]
is therefore compatible with the filtrations
\(\bigl(\omega(r)_n\bigr)\) and
\(\bigl(\mathfrak m(r)_n\bigr)\). Since
\[
  k[[{}^{r}\omega]]
  \big/
  \bigl((x^{p^n})\bigr)_{x\in{}^{r}\omega_n}
\]
is exactly \(B'_r\), and \(\phi_r:B'_r\to B_r\) is induced by
\(\sigma(r)\), the result already proved for finite-height
algebras shows that \(\phi_r\) is an isomorphism.
